%\newpage

\section{Introduction}
%\addcontentsline{toc}{section}{Introduction}

Most of us were subjected to stepping onto the playground in primary school to experience a cornerstone of social organization. This is the idea of \textit{status hierarchies}, which are ``characterized by a rank ordering of individuals or groups according to the amount of respect accorded by others'' \citep{Magee-2008-SocialHierarchy}. Status hierarchies are a basic, yet universal aspect of social organization that extend from humans to non-human animals \citep{Ridgeway-2019-Status, Sapolsky-2004-Social, Halevy-2011-Functional}. Although social scientists have long been engrossed in how status hierarchies develop in human and non-human social groups, an open question remains addressing whether status hierarchies might also emerge among artificial intelligence (AI) systems. Because AI systems, specifically language models, are predominately trained on human data replete with language patterns, preferences, and biases, they may reproduce human social dynamics, including the emergence of status hierarchies. This thesis addresses this by asking: \textit{Do language models form status hierarchies, and if so, through what mechanisms?}

To address this question, I adapt an experimental design from \citet{Berger-1972-Status} that tests whether humans form expectation states, or internalized beliefs about relative competence and influence, based on status characteristics, meaning socially recognized attributes that shape evaluations and interactions within groups, often producing hierarchical dynamics (e.g., job title, educational background, gender, income). In this design, participants completed a subjective decision-making task, made initial choices, then saw their partner's choice. The dependent variable was the percentage of participants sticking to their initial choice or deferring to their partner's choice.

I replicate this core logic with language models. Rather than human participants, I use separate instances of language models---that is, independent runs of models within a shared code script, initialized without shared memory or prior context. The task is now divided into four parts: First, model's read identical texts, in this case movie reviews, and independently provide a sentiment rating. Second, models are introduced to each other, having their status characteristics revealed. Third, each model is shown both its own initial rating and its partners' initial rating. Finally, both models are given the opportunity to maintain or revise their initial rating. The dependent variable remains conceptually identical as the percentage of trials in which a model changes its rating in the direction of its partners' initial rating.

\textbf{[WIP]: In this paragraph I will offer initial, high-level results from the experiments.}

One might easily ask why the emergence of status hierarchies in language models would matter at all. Status competition could introduce unintended behaviors where AI systems prioritize their relative standing over collaborative goals or engage in deceptive strategies that undermine reliability. Dynamics like this have already been observed in multi-agent reinforcement learning, where agents develop deceptive strategies to outcompete peers \citep{Leibo-2017-Multi}, and in reward hacking, where systems exploit evaluation mechanisms while subverting intended goals \citep{Amodei-2016-Concrete}. Beyond reliability concerns, models seeking status could amplify existing biases by learning to associate language patterns, credentials, or demographic markers with higher standing, thus perpetuating inequalities in high-stakes domains. This has already seen in AI systems amplifying discriminatory patterns in resume screening \citep{Raghavan-2019-Mitigating}, recidivism prediction \citep{Angwin-2016-MachineBias}, and occupational stereotyping \citep{Kotek-2023-Gender}; status hierarchies could create compounding layers of such bias.

The scalability of AI systems introduces a further dimension of risk: while human status hierarchies form over time, language models interact across millions of contexts simultaneously, enabling status patterns to crystallize and propagate with unprecedented speed. Algorithmic patterns already spread rapidly—recommendation systems entrench filter bubbles across platforms \citep{Pariser-2011-Filter}, trading algorithms synchronize to trigger flash crashes \citep{Kirilenko-2011-Flash}—and distributed language model deployment could institutionalize status dynamics before they become visible. Perhaps most troublingly, status-conscious models may resist interventions that threaten their perceived standing, a form of misalignment where relative position supersedes human benefit. This parallels goal misgeneralization \citep{Langosco-2021-Goal} and specification gaming \citep{Krakovna-2020-Specification, Bostrom-2014-Superintelligence}, where systems pursue unintended objectives correlated with training rewards, suggesting that status-seeking behavior could prove similarly resistant to correction.

This thesis follows the following structure. First, I outline extant literature on status hierarchies across biological and social systems. Then, I outline existing research on emergent behaviors in language models and identify gaps in understanding status hierarchies in AI systems. Following this, I detail my adaption of \citet{Berger-1972-Status}'s classic experimental design, but here, applied to language models across various setups. Next, I continue with experimental results where I examine whether language models exhibit status-seeking behaviors and how these patterns manifest across different models and conditions. Finally, I conclude by discussing the implications of these findings for AI alignment research and consider both risks and opportunities presented by status-aware AI systems.